With two equal minima, Dumortier, Panazzolo and Roussarie unfold a
triple cycle into three hyperbolic cycles inside an annulus whose
boundary components are crossed in opposite radial directions; because
the total multiplicity in the annulus is even, a fourth cycle is forced
\cite[Theorem~1.3 and p.~1898]{DPR}.

In our multiscale construction the initial trinomial plays a different
role: it must simultaneously supply bounded ordinary brackets and
opposite-sign incoming graphs from large radius. We therefore choose its
lowest odd coefficient so that the limiting lower-axis intercepts agree.
Two strict bounded brackets survive at that coefficient. The curvatures
at the equal-depth minima fix the incoming sign at one negative height,
and the leading odd term gives the opposite sign at a greater depth.

\begin{lemma}[Balanced three-power compact seed]\label{lem:compact-seed}
Let $m_0\ge 5$ be an odd integer and put $d_0=2m_0$.
\begin{enumerate}
  \item[\textup{(i)}] There exists a unique real root $v>1$ of
  \begin{equation}\label{eq:seed-v-poly}
    v^{3m_0-1} - (2m_0-1)(v^{2m_0} + v^{m_0-1}) - 2(m_0-1)(v^{2m_0-1} + v^{m_0}) + 1 = 0.
  \end{equation}
  With $b_0 = \frac{2m_0(v^{m_0}+1)}{(2m_0-1)(v^{m_0-1}-1)}>0$ and $a_0 = 2 + (2-1/m_0)b_0>0$, the
  half-profiles $q_\pm(r)=r^{2m_0}\pm(b_0 r^{2m_0-1}-a_0 r^{m_0})$ on $[0,\infty)$ attain equal global
  minima $q_+(1)=q_-(v)=\mu_0 := -1-(1-1/m_0)b_0<0$ at $\rho_+=1$ and $\rho_-=v$, while $q_-$ has a
  single inner local maximum of height $H_b>0$ in $(0,v)$ and the only positive critical points of
  $q_-$ are this maximum and $v$.
  \item[\textup{(ii)}] Suppose that for a given odd $m_0\ge 5$ there exists a height $h_1>H_b$ at
  which the outer roots $t_\sigma(h_1)>\rho_\sigma$ of $q_\sigma(t_\sigma)=h_1$ satisfy
  $S_+(h_1)<S_-(h_1)$, where $S_\sigma(h)=\int_{\rho_\sigma}^{t_\sigma(h)} q_\sigma'(r)^2\frac{\dd
  r}{r}$, together with the curvature inequality $q_-''(v) > v^2 q_+''(1)$. Then there exist a
  coefficient $b>0$, a compact interval $I_\lambda\subset(-\infty,0)$ containing a unique
  tail-balancing value $\lambda_0\in\operatorname{int}(I_\lambda)$ for
  $P_{\mathrm{seed},\lambda}(x)=x^{d_0}+bx^{d_0-1}+\lambda x^{m_0}$, two disjoint compact positive
  intervals $B_1,B_2\subset(0,\infty)$, and two negative heights
  $h_{\mathrm{seed},2}<h_{\mathrm{seed},1}<0$ such that:
  \begin{itemize}
    \item
    $\Delta(P_{\mathrm{seed},\lambda})=\ell_+(P_{\mathrm{seed},\lambda})-\ell_-(P_{\mathrm{seed},\lambda})$
    has opposite strict signs at the endpoints of $I_\lambda$ and vanishes at $\lambda_0$;
    \item for every $\lambda\in I_\lambda$,
    $\oscD_\lambda(\beta)=b_{+,\lambda}(\beta)-b_{-,\lambda}(\beta)$ has opposite strict signs at the
    endpoints of $B_1$ and of $B_2$;
    \item at $\lambda=\lambda_0$, the inward-graph difference $G_+(h)-G_-(h)$ of
    $P_{\mathrm{seed},\lambda_0}$ is strictly negative at $h=h_{\mathrm{seed},1}$ and strictly
    positive at $h=h_{\mathrm{seed},2}$.
  \end{itemize}
  \item[\textup{(iii)}] The hypotheses of \textup{(ii)} hold with $h_1=d_0$ for all sufficiently large
  odd $m_0$, and explicitly at $m_0=33153$ ($d_0=66306$); they also hold at $m_0=5$ ($d_0=10$) with
  $h_1=20$.
\end{enumerate}
\end{lemma}

\begin{proof}
(i) The polynomial in \eqref{eq:seed-v-poly} has two coefficient sign changes and takes values $+1$ at
$v=0$, $-8(m_0-1)<0$ at $v=1$, and $+\infty$ as $v\to+\infty$, so Descartes' rule gives a unique root
$v\in(1,\infty)$. The reduced derivatives $q_\sigma'(r)/r^{m_0-1}=2m_0 r^{m_0}+\sigma[(2m_0-1)b_0
r^{m_0-1}-m_0 a_0]$ vanish at $r=1$ ($\sigma=+$) and $r=v$ ($\sigma=-$) by definition of $a_0,b_0$. On
$(0,\infty)$, $q_+'(r)/r^{m_0-1}$ increases strictly, while $q_-'(r)/r^{m_0-1}$ decreases on
$(0,\frac{(2m_0-1)(m_0-1)}{2m_0^2}b_0)$ and increases thereafter. Substituting $a_0,b_0$ into
$q_+(1)=q_-(v)$ recovers \eqref{eq:seed-v-poly}; since $q_-(v)=\mu_0<0$ and $q_-(0)=0$, $v$ is the
outer root of $q_-'$ and the earlier positive root is the unique local maximum of height $H_b>0$.

(ii) Assume $S_+(h_1)<S_-(h_1)$ at some $h_1>H_b$ and $q_-''(v)>v^2 q_+''(1)$. As $h\to+\infty$,
$t_\sigma(h)=h^{1/d_0}-\sigma b_0/d_0+o(1)$, and integrating $q_\sigma'(r)^2/r$ termwise gives
$S_+(h)-S_-(h)=\frac{2b_0 d_0}{2d_0-3}h^{(2d_0-3)/d_0}+o(h^{(2d_0-3)/d_0})>0$; fix $h_2>h_1$ with
$S_+(h_2)>S_-(h_2)$. Consider $F_{C,a}(x)=C(x^{d_0}+b_0 x^{d_0-1}-a x^{m_0})$ as $C\to+\infty$ and $a$
varies near $a_0$. Write $q_{\sigma,a}$ for the profiles with variable $a$ and
$m_\sigma(a)=\min q_{\sigma,a}$. Uniformly near $a_0$,
$\ell_\sigma(F_{C,a})/C\to m_\sigma(a)$. For the upper estimate, take a
fixed launch level $Ch$ above all earlier maxima and use the running
minimum in \eqref{eq:envelope-bounds}. For the lower estimate start the
scaled inward graph at $m_\sigma(a)-\delta$, with $\delta>0$. Until it
reaches $m_\sigma(a)-\delta/2$ its total increase is at most
\[
 C^{-2}\int_0^\infty
  \frac{r\dd r}{q_{\sigma,a}(r)-m_\sigma(a)+\delta/2}=O_\delta(C^{-2})
\]
uniformly in $a$. Thus for large $C$ this graph stays below every turn.
The derivative of $m_+(a)-m_-(a)$ at $a_0$ is $-1-v^{m_0}$. Continuity
and strict monotonicity of the true tail gap (Lemma~\ref{lem:tail-inward})
give a unique balance $a_C\to a_0$, with common intercept $\ell_C$.

For $h>H_b$ put $\eps=C^{-2}$. At a turn $T$ above all preceding
extrema, the running minimum follows the profile from its global
minimum to the turn. By \eqref{eq:turn-derivs}, \eqref{eq:envelope-bounds}
and integration by parts, $\eps J_i(T)$ differs from
$\int_{\rho_\sigma(a)}^{\sqrt{2T}}q_{\sigma,a}'(r)^2\dd r/r$ by
$O(\sqrt\eps)$ on compact intervals of such turns. Here
$\rho_\sigma(a)$ is the minimizer. The same bounds make the scaled
outgoing action $O(\sqrt\eps)$ since the outgoing envelope is constant.
The actual turn from $Ch$ tends to the outer inverse of $h$.

Integrate $Y'(T)=\eps e^{-J_i(T)}$ from that turn to a fixed large
cutoff. The action is strictly increasing, so the negative scaled logarithm
of the integral tends to its value at the first endpoint. For a lower
bound integrate on an interval of length $\eps$ there. The infinite
remainder has larger exponential cost. Indeed, on a far rising strip
$[B,4B]$ in $u=r^2/2$ take the profile derivative to be at least $b_*>0$.
Choose a turn beyond the strip and an incoming level above it. Their
inward graphs start on opposite sides of the limiting axis intercept.
Scalar uniqueness keeps the critical inward graph between them and
continues it across the strip.
At $3B$ both have gap $O(\sqrt\eps)$ from the profile by the envelope
bound (truncate the incoming graph beyond the strip with tail error
$O(\eps)$). On $[B,2B]$ both gaps are at most $2\eps/b_*$:
a larger gap would increase at speed at least $b_*/2$ up to $3B$.
Subtraction of the graph equations contracts their difference backwards
by at least $\exp[-b_*^2B/(4\eps)]$. Choose $B$ so that this cost
exceeds both test actions. The passage from squared to signed intercepts
has vanishing scaled logarithmic cost since $m_\sigma(a_0)<0$. This proves
\begin{equation}\label{eq:seed-drop-rate}
 -C^{-2}\log\bigl(b_{\sigma,C}(Ch,a_C)-\ell_C\bigr)\longrightarrow S_\sigma(h).
\end{equation}
At exact balance $\ell_{+,C}=\ell_{-,C}=\ell_C$, subtracting $b_{-,C}(Ch)-\ell_C$ from
$b_{+,C}(Ch)-\ell_C$ therefore gives the sign $-\sgn(S_+(h)-S_-(h))$, so $\oscD_{C,a_C}(Ch_1)>0$ and
$\oscD_{C,a_C}(Ch_2)<0$ for all sufficiently large $C$.

To choose the incoming tests, keep $C$ variable. For any fixed $k_0<\mu_0$, the inward solution
$y_{\sigma,C}(r)$ of $y'=r/(Cq_{\sigma,a_C}(r)-y)$ with limit $Ck_0$ satisfies $0\le
Ck_0-y_{\sigma,C}(r)\le C^{-1}\int_r^\infty \frac{s\dd s}{q_{\sigma,a_C}(s)-k_0}=O(C^{-1})$ uniformly
on $[0,\infty)$, so dominated convergence gives
\[
  C\bigl[G_{+,C}(Ck_0) - G_{-,C}(Ck_0)\bigr] \longrightarrow \int_0^\infty \left(\frac{1}{q_-(r)-k_0}
  - \frac{1}{q_+(r)-k_0}\right)r\dd r \qquad (C\to+\infty).
\]
As $k_0\uparrow\mu_0$, Laplace rescaling around the nondegenerate minima $\rho_+=1$ and $\rho_-=v$
gives $\sqrt{\mu_0-k_0}\int_0^\infty \frac{r\dd r}{q_\sigma(r)-k_0}\to
\frac{\pi\sqrt{2}\,\rho_\sigma}{\sqrt{q_\sigma''(\rho_\sigma)}}$. Because
$v/\sqrt{q_-''(v)}<1/\sqrt{q_+''(1)}$, we may fix $k_1<\mu_0$ close to $\mu_0$ so that the limiting
right-hand side is strictly negative. As $K\to+\infty$, substituting $r=K^{1/d_0}t$ gives
$K^{1-1/d_0}\int_0^\infty\bigl((q_-(r)+K)^{-1}-(q_+(r)+K)^{-1}\bigr)r\dd r\to 2b_0\int_0^\infty
t^{d_0}(1+t^{d_0})^{-2}\dd t>0$, so we may fix $k_2=-K<k_1$ where the limiting right-hand side is
strictly positive. Now fix $C>0$ sufficiently large that $\oscD_{C,a_C}(Ch_1)>0$,
$\oscD_{C,a_C}(Ch_2)<0$, $G_{+,C}(Ck_1)-G_{-,C}(Ck_1)<0$, and $G_{+,C}(Ck_2)-G_{-,C}(Ck_2)>0$ hold
simultaneously. At this fixed $C$, near $(0,0)$ one has $b_0 r^{d_0-1}-a_C r^{m_0}<0$.
On sufficiently small half-arcs the linearized rotation gives both
transverse returns inside the region of strict radial order.
Lemma~\ref{lem:variational}(ii) gives $\oscD_{C,a_C}(\beta_0)<0$ at a small $\beta_0\in(0,Ch_1)$.
Splitting the middle point $Ch_1$ into $Ch_1\mp\eta$ ($\eta>0$ small) produces disjoint brackets
$[\beta_0,Ch_1-\eta]$ and $[Ch_1+\eta,Ch_2]$, which persist on a compact interval
$[a_C-\delta_0,a_C+\delta_0]$ with opposite strict signs of $\Delta$ at its endpoints by
Lemma~\ref{lem:tail-inward}(iii). Rescaling $(x,y)=C^{-1/(d_0-1)}(X,Y)$ sets $b=b_0 C^{1/(d_0-1)}>0$,
$\lambda=-a C^{(d_0-m_0)/(d_0-1)}$, and $h_{\mathrm{seed},i}=C^{d_0/(d_0-1)}k_i$, preserving all
assertions.

(iii) Set $t=r^{m_0}$ and $T=v^{m_0}$ as $m_0\to+\infty$. Substituting $T=m_0 y$ into
\eqref{eq:seed-v-poly} and dividing by $m_0^3$ yields $y^3-4y^2+O((\log m_0)/m_0)=0$ locally uniformly
on $y>0$, so $T/m_0\to 4$, $v\to 1$, and $m_0(b_0/v-1)\to 1$. For
$t=m_0y$, use $b_0/r=(b_0/v)(T/t)^{1/m_0}$ to obtain
$b_0/r=1+(1-\log(y/4))/m_0+o(m_0^{-1})$ locally uniformly for $y>0$.
Writing $q_\sigma$ in $t=m_0 y$ and
dividing by $m_0$ shows $m_0^{-1}q_-(r)\to \Phi(y):=y^2(\log(y/4)-1)+4y$ in $C^2$ on compact subsets
of $(0,\infty)$. Since $\Phi'(y)=2y\log(y/4)-y+4$ satisfies $\Phi'(4)=0$ and has a unique earlier zero
$y_b\in(0,4)$ (not equal to $2$ since $\Phi'(2)=2-4\log 2<0$), at $y_b$ we have
$2y_b\log(y_b/4)=y_b-4$, hence $\Phi(y_b)=2y_b-y_b^2/2<2$, proving $H_b/m_0\to\Phi(y_b)<2$, i.e.,
$H_b<2m_0=d_0$ for large $m_0$. In $t=\sqrt{m_0}\,w$, $m_0^{-1}q_+(r)\to 2w^2$, and the exact action
in $t$ is
$m_0\int t^{1-2/m_0}(\partial_t q_\sigma)^2\dd t$.
On the plus branch the integral starts at $w=m_0^{-1/2}$, where
$(\sqrt{m_0}w)^{-2/m_0}\le1$ and the rescaled derivatives are uniformly
bounded up to the upper inverse. Hence $S_+(2m_0)/m_0^3\to \int_0^1 w(4w)^2\dd w=4$, whereas on the
minus side $S_-(2m_0)/m_0^3\to \int_4^{b_*} y\Phi'(y)^2\dd y$ with $\Phi(b_*)=2$.
Here $4<b_*<8$ since $\Phi$ increases beyond $4$ and
$\Phi(4)=0<2<64\log 2-32=\Phi(8)$. Cauchy--Schwarz gives
$\int_4^{b_*}y\Phi'(y)^2\dd y \ge
(\Phi(b_*)-\Phi(4))^2/\int_4^{b_*}y^{-1}\dd y = 4/\log(b_*/4) > 4/\log 2 > 4$. Finally,
$q_+''(1)=2m_0^2+(2m_0-1)(m_0-1)b_0\sim 4m_0^2$, while $q_-''(v)=v^{2m_0-3}[2m_0^2
v-(2m_0-1)(m_0-1)b_0]\sim 16m_0^3$ (reflecting $\Phi''(4)=1$), so $q_-''(v)/(v^2 q_+''(1))\sim
4m_0\to+\infty$.

At the explicit finite exponent $m_0=33153$ ($d_0=66306$), validated 768-bit Arb ball bisection
\cite{Johansson2017} (Section~\ref{sec:finite-inputs}) certifies $0.81<H_b/d_0<0.82$,
$4<S_+(d_0)/m_0^3<5<16<S_-(d_0)/m_0^3<17$, and $q_-''(v)/(v^2 q_+''(1))>10^5$.

For $m_0=5$ the polynomial in \eqref{eq:seed-v-poly} factors as
\[
 (v+1)^4(v^4-v^3+v^2-v+1)\bigl[(v^2-v+1)^3-3v^3\bigr].
\]
The middle factor equals $(v^5+1)/(v+1)$. Thus $v$ is the unique
root above one of $v+v^{-1}=1+\sqrt[3]{3}$. This is the radius in
Chen, Jin, Li and Zhou \cite[(3.14)--(3.16)]{CJLZ} for their linear,
ninth- and tenth-power profile at equal minima. Here the lowest power
is five, $b_0=10(v^5+1)/(9(v^4-1))$, $a_0=(9b_0+10)/5$, and
the action is evaluated for these coefficients.

For a paper calculation let $I_j(N)=[N,N+1]10^{-j}$.
The sextic $(v^2-v+1)^3-3v^3$ has signs $-,+$ at the ends of
$I_{15}(1921942693683930)$. With this interval for $v$, substitute in
$b_0,a_0$ above. The reduced derivative $10r^5-9b_0r^4+5a_0$ has signs
$+,-$ at the ends of $I_9(1461008998)$, on its descending branch
$r<18b_0/25$. This encloses the inner maximum radius. Substitution in
$q_-$ gives $0<H_b<14$. At the two minima,
\[
 130<q_+''(1)=50+36b_0<145,\qquad
 850<q_-''(v)=v^7(50v-36b_0)<1100,\qquad 145v^2<850 .
\]
For $\sigma=\pm1$ an antiderivative of $(q_\sigma')^2/r$ is
\[
 \begin{aligned}
 \Phi_\sigma(r)={}&\tfrac{100}{18}r^{18}
 +\sigma\tfrac{180}{17}b_0r^{17}+\tfrac{81}{16}b_0^2r^{16}\\
 &-\sigma\tfrac{100}{13}a_0r^{13}
 -\tfrac{90}{12}a_0b_0r^{12}+\tfrac{25}{8}a_0^2r^8 .
 \end{aligned}
\]
The first row, with the curvature bounds, proves (ii) at $h_1=20$.
The two higher levels serve the small-support constructions.
The table gives $I_9(N_\sigma)$ for the two outer radii. In each
column $q_\sigma(N_\sigma10^{-9})<h<q_\sigma((N_\sigma+1)10^{-9})$.
The last column bounds
$\Phi_+(t_+)-\Phi_+(1)-\Phi_-(t_-)+\Phi_-(v)$.
\[
\begin{array}{r|rr|c}
 h &N_+&N_-&S_+(h)-S_-(h)\\ \hline
 20&1320472481&2085892924&(-89,-88)\\
 100&1497737756&2206941520&(1293,1295)\\
 200&1594050028&2279074632&(6884,6886)
\end{array}
\]
For rational interval arithmetic, evaluate a positive product or power at its
left and right endpoints, and for a subtraction $I-J$ use
$[\inf I-\sup J,\sup I-\inf J]$. These rules applied term by term give
the displayed inequalities.
\end{proof}
